\chapter{Control editorial de la plantilla}
\textbf{Este anexo es una guía de edición y debe retirarse de la memoria final.} Se oculta al desactivar \texttt{notaseditoriales} en \texttt{config/datos.tex}.

\section{Cumplimiento de formato}
\begin{longtable}{p{4.5cm}p{8.7cm}}
\toprule\textbf{Requisito} & \textbf{Implementación y revisión}\\\midrule\endhead
Papel carta & Clase \texttt{letterpaper}; verificar PDF final.\\
Margen izquierdo 4 cm & Configurado con \texttt{geometry}; otros márgenes 3 cm.\\
Times New Roman & Detección automática o archivos TTF propios. La vista previa puede usar una alternativa; activar modo estricto antes de entregar.\\
Tamaño del cuerpo & Se adopta 12 pt; la norma no fija explícitamente un tamaño de cuerpo.\\
Interlineado & 1,5 en cuerpo; ajuste entre párrafos para doble espacio; citas a espacio simple.\\
Cubierta y portada & Dos páginas; tamaños 18/16/14/20/14/12/12/14 pt y escudo de 33 mm.\\
Portada & Negrita, correferente y advertencia institucional.\\
Paginación & Inferior derecha. Preliminares e introducción en romanos; desarrollo en arábigos.\\
Conclusiones & Título sin numeración; incluidas en el índice.\\
Figuras y tablas & Numeración por capítulo; título de figura debajo y de tabla encima.\\
Ecuaciones & Numeración por capítulo a la derecha entre paréntesis.\\
Anexos & Letras mayúsculas mediante \texttt{appendix}.\\
Referencias & BibTeX numérico en orden de citación; revisar metadatos conforme al anexo C de la norma.\\
\bottomrule
\end{longtable}

\section{Cuestiones técnicas pendientes}
\begin{enumerate}
\item Congelar la arquitectura y orden de estados del cartpole final; diferenciarlo del prototipo LQR.
\item Resolver la identidad del modelo UART de radiografías: PTQ interno INT8 con interfaz Float32 o mixta selectiva con interfaz INT8.
\item Vincular el ELF del MPC HIL con su perfil; resolver la discrepancia CPU frente a partición CPU/NPU.
\item Revisar la trayectoria HIL completa antes de afirmar estabilidad sostenida desde 4 s.
\item Separar tamaños de archivos, Flash compilada y SRAM máxima.
\item Completar resultados faltantes de clasificación, recursos y energía.
\item Verificar límites del control y criterio de éxito en el script final.
\item Cerrar resumen y traducción, objetivos cumplidos y selección definitiva.
\end{enumerate}

\section{Identidad de fuentes}
Los informes tempranos se usan como antecedentes del proceso y no como evidencia definitiva de la ejecución final. El documento \path{paper_MPC_arquitecturaModificada.pdf} corresponde a control de invernaderos; \texttt{MixPresition.pdf} corresponde a MPC explícito robusto. El artículo de origen del cartpole es \termino{Parameter-Adaptive Approximate MPC: Tuning Neural-Network Controllers without Retraining}.

El DOI de ejemplo que aparece en la copia preliminar de \texttt{MixPresition.pdf} no se utilizó. La bibliografía registra el DOI publicado 10.1145/3358223. Los informes internos incluyen las fechas que se pudieron identificar; no se inventó editorial, DOI o fecha para campos ausentes.

\section{Antes de entregar}
\begin{enumerate}
\item Confirmar título aprobado, denominación profesional y mes del examen.
\item Instalar Times New Roman y activar \texttt{fuenteestricta}.
\item Sustituir todas las figuras pendientes y resolver las celdas vacías relevantes.
\item Completar referencias y revisar citas contra sus fuentes.
\item Desactivar notas editoriales y retirar instrucciones que no pertenecen al informe.
\item Compilar y revisar índices, márgenes, paginación y desbordamientos.
\end{enumerate}
